Partial actions of groups on C*-algebras (initially defined for $\mathbb{Z}$ in \cite{MR1276163} and for general discrete groups in \cite{MR1331978}), are a powerful tool in the study of C*-algebras associated to dynamical systems (\cite{MR3043025,MR3539347,MR1703078,MR3699170,MR3210558,MR3600124}). In \cite{MR2115083} Dokuchaev and Exel introduced partial group actions in a purely algebraic context, and although this algebraic theory is not yet as well-developed as its C*-algebraic counterpart, it has attracted interest from researchers in the area due to important classes of algebras, such as graph and ultragraph Leavitt path algebras, being described as (partial) crossed products (\cite{MR3196063, 1706.03628v2}).

Groupoids are also being extensively used in order to study and classify similar classes of algebras (\cite{MR3274831}), since every partial group action on a topological space induces a transformation groupoid (\cite{MR2045419}), and every étale groupoids can be obtained in this manner (\cite[Proposition 5.4]{MR2419901}). In fact, a categorical relationship between the representation theories of groupoids and of inverse semigroups is described in \cite{MR2969047}.

In the C*-algebraic case, the algebras of transformation groupoids coincide with crossed products of associated partial actions (\cite{MR2045419}). However, in the algebraic case only partial results of this nature had been obtained (\cite{MR3743184, arxiv1710.09723v1}). 

In this chapter we will study partial actions of inverse semigroups, as defined in \cite{MR3231479}. They are a common generalization of both partial actions of groups and actions of inverse semigroups. We will be mostly interested in partial actions on topological spaces and algebras, although the same theory can be extended to other classes of algebraic or geometric spaces. These allows us to construct \emph{groupoids of germs} as in \cite{MR2419901} and crossed products associated to partial actions of inverse semigroups. These constructions generalize transformation groupoids and the usual crossed products by groups, respectively.

In this general setting, we prove that the Steinberg algebra of a Hausdorff groupoid of germs is (naturally) isomorphic to a crossed product algebra for partial actions of inverse semigroups (Theorem \ref{theoremsteinbergiscrossed}). This is a generalization of the aforementioned results of \cite{MR3743184, arxiv1710.09723v1}, and also an algebraic analogue of Abadie's result in \cite{MR2045419} in the general context of inverse semigroups.

Moreover, there are canonical constructions which allow us to build new actions from partial actions of inverse semigroups (\cite{MR2419901,MR3231479}), and so we study the isomorphism of the associated groupoids of germs and crossed product algebras. In particular, we obtain algebraic analogues of some of the results of \cite{MR3231226}.

%In fact the results of \cite{MR2419901} prove that partial group actions can be regarded as actions of inverse semigroups, which were already considered in \cite{MR1456588} and can be used, for example, to describe groupoid C*-algebras as crossed products by inverse semigroups \cite[Theorem~3.3.1]{MR1724106} . Although these approaches are similar in some respects, each of them has its advantages and drawbacks -- for example, actions of inverse semigroups respect the operation completely, whereas groups have, overall, a better algebraic structure than general inverse semigroups.

%The Hausdorff assumption we make on the groupoid of germs has been necessary throughout the recent papers in this direction \cite{MR3539347, Steinberg2017, Cordeiro2017}, and it is always satisfied by semigroups which are weak semilattices, as long as we restrict ourselves to ample actions. Even more strongly, all transformation groupoids of partial group actions on Hausdorff spaces are always Hausdorff, and the same will also be true for all $E$-unitary inverse semigroups (Proposition~\ref{prop:eunitarygroupoidofgerms}). These are the sub-inverse semigroups of semidirect products of lattices by groups (\cite[Theorem~7.1.5]{MR1694900}).

Orbit equivalence and full groups for actions of $\mathbb{Z}$ were initially studied by Giordano, Putnam and Skau (\cite{MR1363826, MR1710743,MR2054051}), and by Li in \cite{MR3789176,MR3694599} for partial actions of discrete groups. The notion of continuous orbit equivalence can be naturally extended to partial actions of inverse semigroups. We also introduce and study a natural notion of freeness for partial inverse semigroup actions (which is more restrictive than the one in \cite{MR3448414}). This notion corresponds to the associated groupoids of germs being topologically principal. We also prove that two ample, topologically free partial inverse semigroup actions are continuously orbit equivalent if and only if their corresponding groupoids of germs are isomorphic.% It is important to note that the semigroups considered do not need to be isomorphic, since continuous orbit equivalence deals mostly with the dynamics of the base space inherited from the partial action (see also Matui 
\nocite{MR2876963,MR3377390}%MATUI
\nocite{MR2861532}%Medynets